% Options for packages loaded elsewhere
\PassOptionsToPackage{unicode}{hyperref}
\PassOptionsToPackage{hyphens}{url}
\documentclass[
]{article}
\usepackage{xcolor}
\usepackage{amsmath,amssymb}
\setcounter{secnumdepth}{-\maxdimen} % remove section numbering
\usepackage{iftex}
\ifPDFTeX
  \usepackage[T1]{fontenc}
  \usepackage[utf8]{inputenc}
  \usepackage{textcomp} % provide euro and other symbols
\else % if luatex or xetex
  \usepackage{unicode-math} % this also loads fontspec
  \defaultfontfeatures{Scale=MatchLowercase}
  \defaultfontfeatures[\rmfamily]{Ligatures=TeX,Scale=1}
\fi
\usepackage{lmodern}
\ifPDFTeX\else
  % xetex/luatex font selection
\fi
% Use upquote if available, for straight quotes in verbatim environments
\IfFileExists{upquote.sty}{\usepackage{upquote}}{}
\IfFileExists{microtype.sty}{% use microtype if available
  \usepackage[]{microtype}
  \UseMicrotypeSet[protrusion]{basicmath} % disable protrusion for tt fonts
}{}
\makeatletter
\@ifundefined{KOMAClassName}{% if non-KOMA class
  \IfFileExists{parskip.sty}{%
    \usepackage{parskip}
  }{% else
    \setlength{\parindent}{0pt}
    \setlength{\parskip}{6pt plus 2pt minus 1pt}}
}{% if KOMA class
  \KOMAoptions{parskip=half}}
\makeatother
\usepackage{longtable,booktabs,array}
\newcounter{none} % for unnumbered tables
\usepackage{calc} % for calculating minipage widths
% Correct order of tables after \paragraph or \subparagraph
\usepackage{etoolbox}
\makeatletter
\patchcmd\longtable{\par}{\if@noskipsec\mbox{}\fi\par}{}{}
\makeatother
% Allow footnotes in longtable head/foot
\IfFileExists{footnotehyper.sty}{\usepackage{footnotehyper}}{\usepackage{footnote}}
\makesavenoteenv{longtable}
\setlength{\emergencystretch}{3em} % prevent overfull lines
\providecommand{\tightlist}{%
  \setlength{\itemsep}{0pt}\setlength{\parskip}{0pt}}
% ===== unicode + compile header for ZIMMERMAN_LAW_OF_GRAVITY =====
% TO COMPILE ON OVERLEAF: upload this .tex + a figures/ folder containing
%   fig1_rar.png  fig2_a0z.png  fig3_btfr.png
% Compile with XeLaTeX (recommended) or pdfLaTeX. Download the PDF and publish
% manually on Zenodo. (Not test-compiled locally; unicode is mapped below.)
\usepackage{amssymb}
\usepackage{newunicodechar}
\providecommand{\textsubscript}[1]{\ensuremath{_{\mbox{\scriptsize #1}}}}
\providecommand{\textonehalf}{\ensuremath{\frac12}}
\newunicodechar{₀}{\ensuremath{_{0}}}
\newunicodechar{₁}{\ensuremath{_{1}}}
\newunicodechar{₃}{\ensuremath{_{3}}}
\newunicodechar{₆}{\ensuremath{_{6}}}
\newunicodechar{ₐ}{\ensuremath{_{a}}}
\newunicodechar{¹}{\ensuremath{^{1}}}
\newunicodechar{²}{\ensuremath{^{2}}}
\newunicodechar{⁰}{\ensuremath{^{0}}}
\newunicodechar{⁴}{\ensuremath{^{4}}}
\newunicodechar{⁵}{\ensuremath{^{5}}}
\newunicodechar{⁻}{\ensuremath{^{-}}}
\newunicodechar{Λ}{\ensuremath{\Lambda}}
\newunicodechar{ρ}{\ensuremath{\rho}}
\newunicodechar{π}{\ensuremath{\pi}}
\newunicodechar{σ}{\ensuremath{\sigma}}
\newunicodechar{Υ}{\ensuremath{\Upsilon}}
\newunicodechar{χ}{\ensuremath{\chi}}
\newunicodechar{Ω}{\ensuremath{\Omega}}
\newunicodechar{Δ}{\ensuremath{\Delta}}
\newunicodechar{γ}{\ensuremath{\gamma}}
\newunicodechar{α}{\ensuremath{\alpha}}
\newunicodechar{µ}{\ensuremath{\mu}}
\newunicodechar{√}{\ensuremath{\surd}}
\newunicodechar{×}{\ensuremath{\times}}
\newunicodechar{−}{\ensuremath{-}}
\newunicodechar{·}{\ensuremath{\cdot}}
\newunicodechar{∝}{\ensuremath{\propto}}
\newunicodechar{≈}{\ensuremath{\approx}}
\newunicodechar{≳}{\ensuremath{\gtrsim}}
\newunicodechar{≥}{\ensuremath{\geq}}
\newunicodechar{⇒}{\ensuremath{\Rightarrow}}
\newunicodechar{→}{\ensuremath{\rightarrow}}
\newunicodechar{↔}{\ensuremath{\leftrightarrow}}
\newunicodechar{⟺}{\ensuremath{\Longleftrightarrow}}
\newunicodechar{⋆}{\ensuremath{\star}}
\newunicodechar{★}{\ensuremath{\star}}
\newunicodechar{✓}{\ensuremath{\checkmark}}
\newunicodechar{✗}{\ensuremath{\times}}
\newunicodechar{◐}{\ensuremath{\circ}}
\newunicodechar{½}{\ensuremath{\frac{1}{2}}}
\newunicodechar{§}{\S{}}
\newunicodechar{Ö}{\"{O}}
\newunicodechar{Ü}{\"{U}}
\newunicodechar{é}{\'{e}}
\newunicodechar{³}{\ensuremath{^{3}}}
\newunicodechar{ä}{\"{a}}
\newunicodechar{Ȳ}{\ensuremath{\bar{Y}}}
\newunicodechar{δ}{\ensuremath{\delta}}
\newunicodechar{θ}{\ensuremath{\theta}}
\newunicodechar{κ}{\ensuremath{\kappa}}
\newunicodechar{λ}{\ensuremath{\lambda}}
\newunicodechar{μ}{\ensuremath{\mu}}
\newunicodechar{φ}{\ensuremath{\varphi}}
\newunicodechar{–}{\textendash{}}
\newunicodechar{—}{\textemdash{}}
\newunicodechar{…}{\ldots{}}
\newunicodechar{′}{'}
\newunicodechar{⁶}{\ensuremath{^{6}}}
\newunicodechar{⁷}{\ensuremath{^{7}}}
\newunicodechar{₂}{\ensuremath{_{2}}}
\newunicodechar{₄}{\ensuremath{_{4}}}
\newunicodechar{ℏ}{\ensuremath{\hbar}}
\newunicodechar{∇}{\ensuremath{\nabla}}
\newunicodechar{≔}{\ensuremath{:=}}
\newunicodechar{≤}{\ensuremath{\leq}}
\newunicodechar{≲}{\ensuremath{\lesssim}}
\newunicodechar{⊙}{\ensuremath{\odot}}
\newunicodechar{Σ}{\ensuremath{\Sigma}}
\newunicodechar{η}{\ensuremath{\eta}}
\newunicodechar{ν}{\ensuremath{\nu}}
\newunicodechar{″}{''}
\newunicodechar{⇔}{\ensuremath{\Leftrightarrow}}
\newunicodechar{±}{\ensuremath{\pm}}
\newunicodechar{¼}{\ensuremath{\tfrac{1}{4}}}
\newunicodechar{β}{\ensuremath{\beta}}
\newunicodechar{∈}{\ensuremath{\in}}
\newunicodechar{∞}{\ensuremath{\infty}}
\newunicodechar{≠}{\ensuremath{\neq}}
% --- v3 erratum build: additional symbol coverage ---
\newunicodechar{ω}{\ensuremath{\omega}}
\newunicodechar{τ}{\ensuremath{\tau}}
\newunicodechar{∫}{\ensuremath{\int}}
\newunicodechar{∂}{\ensuremath{\partial}}
\newunicodechar{≫}{\ensuremath{\gg}}
\newunicodechar{⟨}{\ensuremath{\langle}}
\newunicodechar{⟩}{\ensuremath{\rangle}}
\newunicodechar{⁸}{\ensuremath{^{8}}}
\usepackage{bookmark}
\IfFileExists{xurl.sty}{\usepackage{xurl}}{} % add URL line breaks if available
\urlstyle{same}
\hypersetup{
  hidelinks,
  pdfcreator={LaTeX via pandoc}}

\author{}
\date{}

\begin{document}

\section{Tidal Dwarf Galaxies as a Test of History-Dependent Inertia:
How a Long-Memory Kernel Can Leave Young Dwarfs
Near-Newtonian}\label{tidal-dwarf-galaxies-as-a-test-of-history-dependent-inertia-how-a-long-memory-kernel-can-leave-young-dwarfs-near-newtonian}

\textbf{Carl Zimmerman} Briar Creek Tech (carl@briarcreektech.com)

\begin{center}\rule{0.5\linewidth}{0.5pt}\end{center}

\subsection{Abstract}\label{abstract}

Tidal dwarf galaxies (TDGs) --- condensations formed from tidal debris
pulled out of a progenitor disc during a galaxy interaction --- are the
classic sign-flip test between dark matter and modified dynamics. In
\(\Lambda\)CDM they are born devoid of non-baryonic dark matter and
should be Newtonian (\(M_\mathrm{dyn}/M_\mathrm{bar}\simeq 1\)). In
\emph{instantaneous} modified dynamics --- modified gravity
(AQUAL/QUMOND) or short-memory modified inertia --- they sit deep in the
low-acceleration regime (current baryonic acceleration
\(g_\mathrm{bar}/a_0 = 0.015\)--\(0.073\); Lelli et al.~2015) and should
show a large boost, \(\nu\!\sim\!6\),
\(M_\mathrm{dyn}/M_\mathrm{bar}\sim 4\)--\(8\). The best HI kinematics (Lelli
et al.~2015, six TDGs) instead give
\(M_\mathrm{dyn}/M_\mathrm{bar}=0.9\)--\(1.5\): near-Newtonian. Instantaneous
modified dynamics over-predicts the rotation velocities even with the
external-field effect (EFE), at roughly \(2\)--\(3\sigma\) per NGC 5291
object. This paper reads the same data through a de Sitter--Unruh
\emph{modified-inertia} framework with a \textbf{committed memory
kernel} \(\tau_\mathrm{mem}=2Z/H_\Lambda = 203\,\mathrm{Gyr}\) (the same
kernel used in the relational \(\sigma\)-spread and SHLEM analyses).
Because TDG ages (\(\sim\!0.36\) Gyr) are minute compared with
\(\tau_\mathrm{mem}\), the debris gas has \emph{not} relaxed to its current
deep-regime inertia; its effective inertia should instead be dominated
by the \(\sim\!10\)--\(13\) Gyr it spent at higher acceleration in the
progenitor disc. Averaging the way the framework's own \(\sigma\)-spread
analysis does --- over the convex response \(\nu\), not the acceleration
--- a plausible outer-disc history gives
\(M_\mathrm{dyn}/M_\mathrm{bar}\sim 1.4\)--\(1.9\): this resolves the gross
factor-\(\sim\!4\) instantaneous over-prediction but lands at the
\textbf{high edge of, and mildly above,} the data (over-predicting the
older NGC 7252/VCC within their errors), a consistency rather than a
bullseye. We stress this is a \textbf{plausibility estimate, not a
derived prediction}: the map from the memory-weighted history to an
effective boost is not pinned by the committed formalism (averaging
\(y\) vs the response swings it \(\sim\!20\)--\(30\%\)), so we frame the
result as a hypothesis plus a forecast rather than a firm number. We
report it \emph{together with its cost}: under the long-memory reading
the framework yields the \textbf{same} near-Newtonian TDGs as
\(\Lambda\)CDM, so the classic sign-flip discriminator against dark
matter \textbf{dissolves} for this framework. The instantaneous-MOND
tension is real but re-scoped to modified gravity and short-memory
inertia. We flag three caveats --- (i) Lelli's non-equilibrium escape is
degenerate with the memory reading, (ii) the map from the 203-Gyr
acceleration history to an effective boost \(\nu_\mathrm{eff}\) is not
pinned by the committed formalism (the estimate is kernel-hostage), and
(iii) an earlier pass of this analysis wrongly reported a framework
tension by applying the instantaneous-MOND prediction, which we retract
here. This paper must not be read as ``TDGs show no dark matter'' or as
confirmation of the framework.

\begin{center}\rule{0.5\linewidth}{0.5pt}\end{center}

\subsection{1. Introduction}\label{introduction}

Tidal dwarf galaxies occupy a privileged position in the
dark-matter-versus-modified-dynamics debate. They form out of gas and
stars flung into tidal tails during a major interaction (Bournaud et
al.~2007). In the standard cosmological picture their material comes
from the \emph{baryonic} disc of the progenitor, not from its extended
cold-dark-matter halo; the debris is dynamically cold baryons and any
pre-existing halo cannot be efficiently captured into these small,
late-forming condensations. TDGs should therefore be essentially free of
non-baryonic dark matter and obey Newtonian dynamics on their own
baryonic content, \(M_\mathrm{dyn}/M_\mathrm{bar}\simeq 1\) (Barnes \&
Hernquist 1992; Bournaud et al.~2007).

Modified-dynamics frameworks make the opposite prediction --- \emph{if}
they act instantaneously. The internal accelerations in TDGs are low.
Using the Lelli et al.~(2015) kinematics, the current baryonic
acceleration spans \(g_\mathrm{bar}/a_0 = 0.015\)--\(0.073\), i.e.~deep in
the low-acceleration regime. Any theory whose boost function responds to
the \emph{instantaneous} baryonic acceleration --- modified gravity in
its AQUAL/QUMOND realizations, or a modified-inertia theory with a short
response time --- then demands a large dynamical enhancement,
\(\nu = \sqrt{1+1/y}\to 5.9\) at \(y\equiv g_\mathrm{bar}/a_0 = 0.03\), and
hence \(M_\mathrm{dyn}/M_\mathrm{bar}\sim 4\)--\(8\). This is the sign flip:
dark matter says Newtonian, instantaneous modified dynamics says boosted
(Gentile et al.~2007; Lelli et al.~2015). TDGs are one of the few
systems where the two pictures diverge cleanly \emph{and} the
discriminating variable is a directly observed rotation velocity.

The framework examined here is a de Sitter--Unruh
\textbf{modified-inertia} theory. Inertia is treated as a
\emph{time-nonlocal} response of a body to the de Sitter--Unruh bath it
experiences at its own proper acceleration (Milgrom 1994, 1999 for the
nonlocal-MI programme and the interpolation kernel). Its acceleration
scale is horizon-derived, \(a_0 = cH_\Lambda/Z\), with
\(Z=\sqrt{32\pi/3}=5.789\) posited; this yields a canonical
\(a_0 = 9.355\times10^{-11}\,{\rm m\,s^{-2}}\) and an alternative
footing \(a_0 = 1.1305\times10^{-10}\,{\rm m\,s^{-2}}\), both carried
throughout. The load-bearing feature for the present analysis is not the
value of \(a_0\) but the \emph{memory time}: the framework commits to a
response kernel with

\[
\tau_\mathrm{mem} = \frac{2Z}{H_\Lambda} = 203\ \text{Gyr},
\]

the same kernel that produced the SHLEM null (Zenodo
10.5281/zenodo.21458605) and the relational \(\sigma\)-spread signature
(Zenodo 10.5281/zenodo.21421896). A memory that long has an immediate
consequence for \emph{young} systems, and TDGs are the youngest
galaxy-scale objects with clean kinematics. That is the thesis: on this
framework's own committed kernel, TDGs cannot have relaxed into their
deep-regime inertia state, so the instantaneous-MOND prediction does not
apply to them --- with a cost we state up front.

\begin{center}\rule{0.5\linewidth}{0.5pt}\end{center}

\subsection{2. The data and the instantaneous-modified-dynamics
tension}\label{the-data-and-the-instantaneous-modified-dynamics-tension}

The kinematic data are from Lelli et al.~(2015, A\&A 584, A113 =
arXiv:1509.05404), the current best HI observations of TDGs, superseding
the earlier and higher velocity estimates of Gentile et al.~(2007) that
had reported consistency with modified dynamics. Lelli's re-analysis
finds systematically \emph{lower} circular velocities, which is what
makes the honest test bite. We use six TDGs: the three condensations in
the NGC 5291 ring (N, S, SW), two in NGC 7252 (E, NW), and VCC 2062.

For each object we take the baryonic acceleration \(g_\mathrm{bar}/a_0\),
the observed circular velocity \(V_\mathrm{obs}\pm\delta V\), the baryonic
mass, and the observed \(M_\mathrm{dyn}/M_\mathrm{bar}\), together with
Lelli's own MOND predictions for the isolated case (\(V_\mathrm{ISO}\)) and
with the external-field effect (\(V_\mathrm{EFE}\)), computed at their
\(a_0 = 1.2\times10^{-10}\). In deep-regime modified dynamics the
predicted velocity scales as \(V^4\propto a_0\), so moving to the
framework footings rescales the velocity by
\((a_{0,\rm fw}/a_{0,\rm std})^{1/4}\): a factor \(0.940\) for the
canonical footing and \(0.985\) for the alternative. The lower framework
\(a_0\) therefore \emph{reduces} the predicted velocity by only
\(\sim\!6\%\) --- nowhere near enough to remove the tension.

{\def\LTcaptype{none} % do not increment counter
\begin{longtable}[]{@{}
  >{\raggedright\arraybackslash}p{(\linewidth - 12\tabcolsep) * \real{0.1429}}
  >{\raggedright\arraybackslash}p{(\linewidth - 12\tabcolsep) * \real{0.1429}}
  >{\raggedright\arraybackslash}p{(\linewidth - 12\tabcolsep) * \real{0.1429}}
  >{\raggedright\arraybackslash}p{(\linewidth - 12\tabcolsep) * \real{0.1429}}
  >{\raggedright\arraybackslash}p{(\linewidth - 12\tabcolsep) * \real{0.1429}}
  >{\raggedright\arraybackslash}p{(\linewidth - 12\tabcolsep) * \real{0.1429}}
  >{\raggedright\arraybackslash}p{(\linewidth - 12\tabcolsep) * \real{0.1429}}@{}}
\toprule\noalign{}
\begin{minipage}[b]{\linewidth}\raggedright
TDG
\end{minipage} & \begin{minipage}[b]{\linewidth}\raggedright
\(g_\mathrm{bar}/a_0\)
\end{minipage} & \begin{minipage}[b]{\linewidth}\raggedright
\(V_\mathrm{obs}\) (km/s)
\end{minipage} & \begin{minipage}[b]{\linewidth}\raggedright
\(V_\mathrm{ISO}\) (canon)
\end{minipage} & \begin{minipage}[b]{\linewidth}\raggedright
\(V_\mathrm{EFE}\) (canon)
\end{minipage} & \begin{minipage}[b]{\linewidth}\raggedright
tension (canon, +EFE)
\end{minipage} & \begin{minipage}[b]{\linewidth}\raggedright
\(M_\mathrm{dyn}/M_\mathrm{bar}\)
\end{minipage} \\
\midrule\noalign{}
\endhead
\bottomrule\noalign{}
\endlastfoot
NGC 5291 N & 0.073 & \(45\pm9\) & 72 & 58 & \(+1.5\sigma\) &
\(1.5\pm0.7\) \\
NGC 5291 S & 0.033 & \(35\pm6\) & 71 & 51 & \(+2.6\sigma\) &
\(1.3\pm0.5\) \\
NGC 5291 SW & 0.033 & \(28\pm7\) & 58 & 43 & \(+2.2\sigma\) &
\(1.2\pm0.7\) \\
NGC 7252 E & 0.019 & \(18\pm5\) & 49 & 28 & \(+2.0\sigma\) &
\(0.9\pm0.6\) \\
NGC 7252 NW & 0.015 & \(21\pm6\) & 59 & 39 & \(+2.9\sigma\) &
\(1.0\pm0.6\) \\
VCC 2062 & 0.022 & \(16\pm7\) & 39 & 25 & \(+1.3\sigma\) &
\(1.0\pm0.9\) \\
\end{longtable}
}

\emph{(Velocities on the canonical framework footing
\(a_0=9.355\times10^{-11}\); the alternative footing
\(1.1305\times10^{-10}\) raises each by \(\sim\!5\%\) relative to
canonical and does not change the sign of the tension. \(V_\mathrm{ISO}\),
\(V_\mathrm{EFE}\) from Lelli et al.~2015, rescaled by the deep-regime
\(V^4\propto a_0\) law.)}

The pattern is unambiguous. All six TDGs are deep in the
low-acceleration regime (\(g_\mathrm{bar}/a_0 = 0.015\)--\(0.073\), every
value well below unity), so any instantaneous boost function demands
\(M_\mathrm{dyn}/M_\mathrm{bar}\sim 1/\sqrt{y}\sim 4\)--\(8\). The data show
\(M_\mathrm{dyn}/M_\mathrm{bar}=0.9\)--\(1.5\): near-Newtonian. Equivalently,
the instantaneous prediction over-predicts the rotation velocity, and it
does so \emph{even with} the external-field effect (the EFE columns):
the NGC 5291 trio is over-predicted at \(+1.5/+2.6/+2.2\sigma\), and the
full set at \(+1.5/+2.6/+2.2/+2.0/+2.9/+1.3\sigma\).

A naive quadrature over all six gives \(5.3\sigma\), but this is not the
honest figure: the six objects share systematics --- a common EFE model,
correlated inclination and distance assumptions, and a shared
equilibrium assumption --- so the effective tension is closer to the
per-object \(\sim\!2\)--\(3\sigma\) seen in the NGC 5291 trio. We quote
the honest range: \textbf{instantaneous modified dynamics over-predicts
TDG velocities at \(\sim\!2\)--\(3\sigma\)}, and the EFE is not a
sufficient escape. This is a genuine tension for modified gravity
(AQUAL/QUMOND) and for any short-memory modified-inertia theory; it
flips the Gentile et al.~(2007) consistency, which rested on higher
assumed velocities.

\begin{center}\rule{0.5\linewidth}{0.5pt}\end{center}

\subsection{3. The framework-first re-read: a frozen disc-era
inertia}\label{the-framework-first-re-read-a-frozen-disc-era-inertia}

The tension in Section 2 is a statement about theories that respond to
the \emph{current} acceleration. The framework studied here does not.
Its inertia is a time-nonlocal response with a committed memory time
\(\tau_\mathrm{mem}=2Z/H_\Lambda = 203\,\mathrm{Gyr}\). The relevant
comparison is between this memory time and the age of the material whose
inertia we are trying to predict.

TDGs are young. The NGC 5291 system formed in a near head-on collision
roughly \(360\,\mathrm{Myr}\) ago; the debris in the other systems is of
comparable or somewhat greater but still sub-Gyr age. Taking a
representative TDG age of \(\sim\!0.36\) Gyr, we have

\[
\tau_\mathrm{mem} = 203\ \mathrm{Gyr}\ \gg\ t_\mathrm{TDG}\sim 0.36\ \mathrm{Gyr},
\]

and indeed \(\tau_\mathrm{mem}\) exceeds the age of the universe. On the
framework's own kernel the debris gas has had essentially no time to
relax toward the deep-regime inertia state that its \emph{current} low
acceleration would eventually imply. Its effective inertia is instead
dominated by the acceleration history it accumulated over the
\(\sim\!10\)--\(13\) Gyr it spent as part of the progenitor's disc,
before the interaction expelled it into the tail.

This is the same duty-factor argument used in the SHLEM analysis: what
enters the boost is a memory-weighted average of the acceleration
history, not the instantaneous value, and when the memory window vastly
exceeds the time spent in the current state the average is pinned to the
earlier, higher-acceleration epoch. Two features of the \emph{committed}
formalism make this freeze robust rather than ad hoc. First, the
field-theory closure carries a single scale \(a_0\) and rejects any
orbital-time (dynamical) relaxation corner as a new scale absent from
the action; the 203-Gyr memory is therefore orbit-independent, and the
debris does not re-equilibrate on a galaxy dynamical time (the natural
first objection). Second, the kernel is power-law and non-saturating, so
the \emph{residence-time} weight of the recent low-acceleration epoch is
\(w\sim t_\mathrm{TDG}/t_\mathrm{cosmic}\sim 0.36/13.8 \simeq 2.6\%\) ---
small, but not vanishing.

To turn this into a number we must \emph{assume} a representative
disc-era acceleration. Tidal tails are drawn preferentially from the
\emph{outer} disc, where \(y_\mathrm{hist}\sim 0.5\)--\(2\) (the
least-bound gas skews toward the low-\(y\), higher-boost end).
Crucially, we average the way the framework's \emph{own}
\(\sigma\)-spread analysis does --- over the \textbf{response} \(\nu\)
(a Jensen gap over the convex \(\nu\)), not over the acceleration \(y\)
--- giving
\(\nu_\mathrm{eff}=(1-w)\,\nu(y_\mathrm{hist})+w\,\nu(y_\mathrm{now})\) with
\(\nu(y_\mathrm{now}{=}0.03)=5.9\). This is the framework-faithful
prescription; averaging \(y\) instead (the generic MOND shortcut) yields
a systematically \emph{lower}, more favorable boost, and the
\(\sim\!20\)--\(30\%\) gap between the two is the dominant
kernel-hostage uncertainty (Caveat 2). We emphasize \(y_\mathrm{hist}\) is
an assumption, not a trajectory-integrated result, so the entries below
are a plausibility estimate, not a derived prediction:

{\def\LTcaptype{none} % do not increment counter
\begin{longtable}[]{@{}
  >{\raggedright\arraybackslash}p{(\linewidth - 4\tabcolsep) * \real{0.3333}}
  >{\raggedright\arraybackslash}p{(\linewidth - 4\tabcolsep) * \real{0.3333}}
  >{\raggedright\arraybackslash}p{(\linewidth - 4\tabcolsep) * \real{0.3333}}@{}}
\toprule\noalign{}
\begin{minipage}[b]{\linewidth}\raggedright
outer-disc history
\end{minipage} & \begin{minipage}[b]{\linewidth}\raggedright
response-averaged \(\nu_\mathrm{eff}\)
\end{minipage} & \begin{minipage}[b]{\linewidth}\raggedright
predicted \(M_\mathrm{dyn}/M_\mathrm{bar}\)
\end{minipage} \\
\midrule\noalign{}
\endhead
\bottomrule\noalign{}
\endlastfoot
\(y_\mathrm{hist}\sim 0.5\) & 1.84 & \(\sim1.8\) \\
\(y_\mathrm{hist}\sim 1.0\) & 1.53 & \(\sim1.5\) \\
\(y_\mathrm{hist}\sim 2.0\) & 1.35 & \(\sim1.4\) \\
\end{longtable}
}

For contrast, the instantaneous prediction at the TDGs' \emph{current}
\(y\sim0.03\) is \(\nu=5.9\) --- the source of the Section 2
over-prediction. The framework's long-memory, response-averaged estimate
is instead \(M_\mathrm{dyn}/M_\mathrm{bar}\sim 1.4\)--\(1.9\) (central
\(\sim\!1.5\), higher if the least-bound-gas prior pushes
\(y_\mathrm{hist}\) low). This \textbf{resolves the gross
factor-\(\sim\!4\) instantaneous over-prediction}
(\(\nu\!\sim\!6\!\to\!\sim\!1.5\)) but lands at the \textbf{high edge
of, and mildly above,} the observed range (NGC 5291 trio
\(1.2\)--\(1.5\), NGC 7252 and VCC 2062 \(0.9\)--\(1.0\)): it
over-predicts the older, lower-mass NGC 7252/VCC by a few tenths ---
within their \(\pm0.6\) errors, but a mild over-prediction, not a
bullseye. The consequence is asymmetric across the three pictures, and
it is worth stating precisely what is and is not established. What
follows firmly from the committed kernel is \emph{qualitative}: a
203-Gyr memory forbids a 0.36-Gyr object from having relaxed into its
current-acceleration boost, so the framework does \textbf{not} inherit
the Section 2 over-prediction --- that \(\sim\!2\)--\(3\sigma\) tension
belongs to instantaneous modified gravity and short-memory modified
inertia, which lack the long kernel. What does \textbf{not} follow
firmly is the \emph{quantitative} landing at \(1.4\)--\(1.9\): that
number rides on the assumed \(y_\mathrm{hist}\) and on the averaging rule,
so it is a plausibility estimate broadly consistent with (though at the
high edge of) the data, not a prediction the framework forces. We
therefore advance the near-Newtonian-young-TDG result as a hypothesis
and a forecast (§5), not as a firm prediction.

\begin{center}\rule{0.5\linewidth}{0.5pt}\end{center}

\subsection{4. The honest ledger}\label{the-honest-ledger}

We now record what is gained and what is lost, both ways, with no
manufactured win.

\textbf{Gain.} TDGs become a candidate \emph{history-dependence}
discriminator. Among the three pictures on the table, instantaneous
modified dynamics (modified gravity, short-memory inertia) is the one
disfavored by the Lelli (2015) data at \(\sim\!2\)--\(3\sigma\), while
long-memory modified inertia and \(\Lambda\)CDM both survive. This makes
TDGs a new member of the framework's history-dependence family alongside
the relational \(\sigma\)-spread and SHLEM --- and, unlike those, it is
a system where existing data are at least compatible with the
memory-kernel direction rather than merely permitting the mechanism in
the abstract. The internal-consistency value is real but bounded: what
the committed 203-Gyr kernel forces is the \emph{qualitative} verdict (a
0.36-Gyr object cannot have relaxed into its current-acceleration boost,
so no large deep-regime enhancement is expected); the
\emph{quantitative} agreement with the observed
\(M_\mathrm{dyn}/M_\mathrm{bar}\) additionally requires the assumed outer-disc
history and averaging rule of Caveat 2, and should be read as
compatibility, not as a fit landed without adjustment.

\textbf{Cost (stated prominently).} Under the long-memory reading, the
framework predicts the \emph{same} near-Newtonian TDGs as
\(\Lambda\)CDM. The classic sign-flip discriminator --- the whole reason
TDGs were considered a decisive modified-dynamics test ---
\textbf{dissolves} for this framework. TDGs do \emph{not} distinguish
long-memory modified inertia from dark matter. This paper therefore must
not be read as ``TDGs show no dark matter'' nor as confirmation of the
framework. It shows only that the framework is \emph{consistent} with
the data, at the price of surrendering TDGs as an anti-dark-matter
argument.

\textbf{Caveat 1 --- non-equilibrium degeneracy.} Lelli's own,
well-motivated escape for instantaneous MOND is that TDG orbital times
greatly exceed their ages, so the discs may not be dynamically settled;
an unrelaxed disc under-estimates the circular velocity and hence the
enclosed potential. Lelli explicitly states the dark-matter-free
conclusion holds \emph{only} assuming equilibrium. That same
non-equilibrium correction would rescue instantaneous MOND by raising
the true velocities toward the boosted prediction. Consequently
``long-memory modified inertia'' and ``instantaneous MOND plus
non-equilibrium'' are \textbf{degenerate} on current data: the
history-dependence reading is \emph{consistent} with the data but is not
\emph{uniquely selected} by it.

\textbf{Caveat 2 --- the \(\nu_\mathrm{eff}\) mapping is kernel-hostage.}
The step from a 203-Gyr acceleration-history average to a specific
effective boost \(\nu_\mathrm{eff}\sim1.4\)--\(1.9\) is \emph{not} pinned
by the committed formalism, in two ways. (a) It assumes a plausible
outer-disc history \(y_\mathrm{hist}\sim 0.5\)--\(2\) rather than a
trajectory-integrated value; the least-bound-gas prior favours the
low-\(y\), higher-\(\nu_\mathrm{eff}\) end. (b) The averaging \emph{rule}
itself is unpinned: averaging the response \(\nu\) (the framework's own
\(\sigma\)-spread prescription, adopted here) gives
\(\nu_\mathrm{eff}\sim1.4\)--\(1.9\), whereas averaging the acceleration
\(y\) (the generic shortcut) gives a lower, more favorable
\(\sim1.2\)--\(1.7\) --- a \(\sim\!20\)--\(30\%\) swing that is the
dominant systematic. The \(\nu_\mathrm{eff}\) figure is a plausibility
estimate, not a derived prediction.

\textbf{Caveat 3 --- self-correction on the record.} The first pass of
this analysis reported a \(\sim\!2\)--\(3\sigma\) tension \emph{against
the framework}, obtained by applying the instantaneous-MOND
(standard-lens) prediction to a framework that commits to a 203-Gyr
kernel. That was a category error: the framework is not an instantaneous
theory, and its committed kernel forbids the deep-regime boost for
objects this young. We retract the adverse framework verdict and
re-scope the tension to instantaneous modified gravity and short-memory
inertia. This correction is preserved in the reference script (its
CORRECTION block) and is stated here transparently rather than quietly
overwritten.

The net entry: the framework \textbf{passes on its own premises}; the
adverse verdict of the first pass is retracted \emph{as a framework
claim} and re-scoped; discriminating power against \(\Lambda\)CDM is
lost; a candidate modified-inertia-versus-modified-gravity lever is
gained, hostage to equilibrium systematics and to the kernel mapping.

\begin{center}\rule{0.5\linewidth}{0.5pt}\end{center}

\subsection{5. Relation to the history-dependence family, and what would
break the
degeneracy}\label{relation-to-the-history-dependence-family-and-what-would-break-the-degeneracy}

The TDG result is not an isolated re-interpretation. The 203-Gyr kernel
is a \emph{single} commitment that has already been carried into two
other analyses: the relational \(\sigma\)-spread signature (Zenodo
10.5281/zenodo.21421896), which introduced the memory time and its
history-dependence class, and the SHLEM null (Zenodo
10.5281/zenodo.21458605), which shares the identical kernel and
duty-factor logic. The formal basis for the kernel is set out in the
modified-inertia field-theory results (Zenodo 10.5281/zenodo.21403470).
TDGs add a third system in which the same long memory has an observable
consequence --- here, freezing young debris into a near-Newtonian state
--- and the consequence is data-facing rather than a forecast.

What the family shares is also its present limitation: every member is
degenerate with a more conventional escape on current data
(non-equilibrium here; astrophysical scatter and modelling assumptions
elsewhere). Breaking the TDG degeneracy specifically requires one of the
following.

\begin{itemize}
\tightlist
\item
  \textbf{Equilibrium-robust kinematics.} The single biggest confound is
  whether the observed \(V_\mathrm{circ}\) reflects a settled potential.
  Stellar and gas kinematics from ELT-class instruments or from
  wide-field integral-field spectroscopy (MUSE and successors), combined
  with dynamical-state indicators, could establish whether the
  near-Newtonian \(M_\mathrm{dyn}/M_\mathrm{bar}\) is intrinsic or an artifact
  of unsettled discs. If the discs are demonstrably settled, the
  non-equilibrium escape for instantaneous MOND closes, and the
  disfavouring of instantaneous modified dynamics --- not of the
  long-memory framework --- sharpens.
\item
  \textbf{An age-versus-discrepancy gradient across TDGs.} The
  distinctive, MOND-impossible prediction of the long-memory reading is
  that the \emph{degree} of near-Newtonian behaviour should correlate
  with age. Very young debris (\(t\ll\tau_\mathrm{relax}\)) should be most
  Newtonian (inertia frozen at the disc-era value); progressively older
  TDGs should drift toward the deep-regime boost as their inertia begins
  to relax to the current low acceleration. A measured gradient of
  \(M_\mathrm{dyn}/M_\mathrm{bar}\) with TDG age --- flat and near unity for
  the youngest, rising for the oldest --- would be a genuine
  history-dependence signature that neither \(\Lambda\)CDM (which
  predicts Newtonian at all ages) nor instantaneous MOND (which predicts
  a large boost at all ages) can produce. \textbf{A candid warning on
  this forecast: within the present six-object sample the sign runs the
  \emph{wrong} way} --- the more evolved NGC 7252 debris shows
  \emph{lower} \(M_\mathrm{dyn}/M_\mathrm{bar}\) (\(0.9\)--\(1.0\)) than the
  younger NGC 5291 (\(1.2\)--\(1.5\)), the opposite of an
  older-is-more-boosted trend. This is within the \(\pm0.6\) errors and
  confounded by mass and geometry, but it is not supporting evidence,
  and an honest test must confront it. Assembling a TDG sample spanning
  a range of ages, with uniform kinematics, is the observational
  programme that would turn the present consistency into a test --- and
  could equally kill the reading.
\end{itemize}

Neither route is available on the Lelli (2015) sample alone, which is
why the current status is \emph{consistency}, not confirmation.

\begin{center}\rule{0.5\linewidth}{0.5pt}\end{center}

\subsection{6. Conclusion}\label{conclusion}

Tidal dwarf galaxies are the textbook sign-flip test between dark matter
and modified dynamics, and on the best available kinematics (Lelli et
al.~2015) they are near-Newtonian:
\(M_\mathrm{dyn}/M_\mathrm{bar}=0.9\)--\(1.5\), with instantaneous modified
dynamics over-predicting the rotation velocities at
\(\sim\!2\)--\(3\sigma\) per NGC 5291 object even with the
external-field effect, and the lower framework \(a_0\) helping by only
\(\sim\!6\%\). Read through the standard instantaneous lens this is a
tension. Read through the de Sitter--Unruh modified-inertia framework's
\emph{own} committed 203-Gyr memory kernel, it is not: TDGs are far too
young (\(\sim\!0.36\) Gyr) to have relaxed out of their inherited
disc-era inertia, so the framework does not inherit the deep-regime
boost, and on the framework's own response-averaging rule it lands at
\(M_\mathrm{dyn}/M_\mathrm{bar}\sim1.4\)--\(1.9\) --- resolving the gross
factor-\(\sim\!4\) over-prediction but sitting at the high edge of /
mildly above the data, a plausibility estimate rather than a firm
prediction, since the map from the memory average to the boost is not
pinned by the committed formalism. The tension is real but belongs to
instantaneous modified gravity and short-memory inertia, not to this
long-memory framework --- a correction we make explicitly, having
initially mis-assigned it.

The cost is equally explicit: the framework thereby predicts the same
near-Newtonian TDGs as \(\Lambda\)CDM, so TDGs stop being a
discriminator against dark matter for this theory. The gain --- TDGs as
a candidate history-dependence lever, disfavouring instantaneous
modified dynamics --- is degenerate with Lelli's non-equilibrium escape
and hostage to a \(\nu_\mathrm{eff}\) mapping that the committed formalism
does not pin down. Whether the lever becomes a test depends on
equilibrium-robust kinematics and, most decisively, on the search for an
age-versus-discrepancy gradient across a uniform TDG sample. The values
of \(a_0\), of \(Z\), and of the map from the 203-Gyr history average to
\(\nu_\mathrm{eff}\) remain posited or estimated, not derived. This paper
reports a consistency and a lost discriminator, not a confirmation.

\begin{center}\rule{0.5\linewidth}{0.5pt}\end{center}

\subsection{References}\label{references}

\begin{itemize}
\tightlist
\item
  Barnes, J. E., \& Hernquist, L. 1992, \emph{Nature}, 360, 715 ---
  tidal-tail formation of dwarf condensations.
\item
  Bournaud, F., et al.~2007, \emph{Science}, 316, 1166 --- missing mass
  / kinematics in tidal dwarf galaxies.
\item
  Gentile, G., Famaey, B., Combes, F., et al.~2007, \emph{A\&A}, 472,
  L25 --- TDGs as a MOND test (earlier, higher assumed velocities;
  reported consistency).
\item
  Lelli, F., Duc, P.-A., Brinks, E., et al.~2015, \emph{A\&A}, 584, A113
  (arXiv:1509.05404) --- best HI kinematics of six TDGs;
  \(g_\mathrm{bar}/a_0\), \(V_\mathrm{ISO}\), \(V_\mathrm{EFE}\),
  \(M_\mathrm{dyn}/M_\mathrm{bar}\); non-equilibrium caution.
\item
  Milgrom, M. 1983, \emph{ApJ}, 270, 365 --- the low-acceleration scale
  \(a_0\).
\item
  Milgrom, M. 1994, \emph{Ann. Phys.}, 229, 384 --- nonlocal
  modified-inertia programme.
\item
  Milgrom, M. 1999, \emph{Phys. Lett. A}, 253, 273 --- the interpolation
  kernel \(\nu=\sqrt{1+1/y}\) (Eq. 9).
\item
  Kroupa, P. 2012, \emph{PASA}, 29, 395 --- dynamical-state cautions for
  tidal dwarfs.
\item
  Zimmerman, C. --- relational \(\sigma\)-spread signature and the
  203-Gyr commitment, Zenodo 10.5281/zenodo.21421896.
\item
  Zimmerman, C. --- SHLEM null (same kernel), Zenodo
  10.5281/zenodo.21458605.
\item
  Zimmerman, C. --- modified-inertia field-theory results, Zenodo
  10.5281/zenodo.21403470.
\end{itemize}

\emph{Novelty note: a literature search found no prior long-memory
modified-inertia reading of the TDG data; ``no duplicate found'' is not
``none exists.'' This work is scoped to the \(a_0\) reframing and the
history-dependence family, consistent with the author's 2026-06-23
retraction of broader claims.}

\end{document}
